\section{What the exact-half forest theorem gives for edge frequencies}
\label{sec:new-edge-half-reductions}

Write $c(G)$ for the number of cycles and $c_e(G)$ for the number containing $e$.
The edge-half conjecture asks for $2c_e(G)\le c(G)$ at some edge of each connected simple cubic graph, apart from $K_4$ and $K_{3,3}$.
Its denominator is the actual cycle count, not the size of the cycle space.

\begin{proposition}[High-girth graphs at or above half density]
\label{prop:new-edge-half-dense-high-girth}
There are absolute integers $D,R$ with the following property.
Let $G$ be a connected cubic graph of rank $r\ge R$ and girth at least $D$.
Every edge of every shortest cycle satisfies
\[
 c_e(G)\le2^{r-2}.
\]
If also $c(G)\ge2^{r-1}$, all those edges satisfy the edge-half inequality.
\end{proposition}

\begin{proof}
The cubic Moore bound gives a shortest cycle $W$ with at most $3\log_2r$ vertices at these ranks.
The established exact-half forest theorem applies to the connected marker $P=V(W)$.
For each cycle $C$ containing a selected edge $e\in E(W)$, the two even states
\[
 E(C),\qquad E(C)\mathbin\triangle E(W)
\]
restrict to forests outside $P$.
All $2c_e(G)$ states are distinct.
Within each group this follows from injectivity, and the two groups differ in their $e$-coordinate.
The forest theorem therefore gives $2c_e(G)\le2^{r-1}$.
Comparison with $c(G)\ge2^{r-1}$ proves the relative assertion.
\end{proof}

\begin{corollary}[A below-half reduction without small cyclic cuts]
\label{cor:new-edge-half-below-density}
For all sufficiently large ranks, a connected simple cubic graph with no cyclic cut of size at most four and with $c(G)\ge2^{r-1}$ satisfies the edge-half conjecture.
\end{corollary}

\begin{proof}
At girth at least $D$, apply Proposition~\ref{prop:new-edge-half-dense-high-girth}.
For the finitely many smaller girths, the established bounded-girth deficit gives $c(G)<2^{r-1}$ at sufficiently large rank.
This contradicts the hypothesis in the remaining cases.
\end{proof}

This reduction leaves genuine relative counting questions when $c(G)<2^{r-1}$.
It also leaves all finite ranks below the stated threshold.
There is no constant in the edge-half conjecture that can absorb those ranks.

\begin{proposition}[The scope of the weighted tripole statement]
\label{prop:new-tripole-near-half}
Assume that every connected simple normalized three-pole $N$, for all weights $w_{ij}\ge0$, has an edge satisfying
\[
 2c_e(N)+2\sum_{i<j}w_{ij}p_{ij}(e)
 \le c(N)+1+\sum_{i<j}w_{ij}p_{ij}.
\]
Then every three-connected simple cubic graph $G$ has an edge with
\[
 2c_e(G)\le c(G)+1.
\]
In particular, the edge-half conclusion holds whenever $c(G)$ is even.
\end{proposition}

\begin{proof}
Delete any vertex $v$.
The resulting graph $N$ is connected, and its three distinct neighbors become the degree-two terminals.
Every cycle through $v$ is exactly one terminal path in $N$, completed by its two edges to $v$.
Consequently
\[
 c(G)=c(N)+\sum_{i<j}p_{ij},\qquad
 c_e(G)=c_e(N)+\sum_{i<j}p_{ij}(e)
\]
for each edge of $N$.
Apply the assumed statement with all weights equal to one.
If $c(G)$ is even, the integer $2c_e(G)$ cannot equal the odd integer $c(G)+1$.
Thus it is at most $c(G)$.
\end{proof}

The weighted statement is still a research target.
The proposition records a conditional implication, not a proof of that target.
For $K_4$ and $K_{3,3}$, every edge attains $(c(G)+1)/2$, so the extra unit cannot be omitted from this general implication.
